% =============================================================================
%  Robótica (MEEC, IST) — Lab 1: simulação do swing de golfe em MuJoCo
%  Compilar com:  latexmk -pdf main.tex
%  Figuras em Photos/ ; bibliografia em refs.bib (bibtex, estilo IEEEtran)
% =============================================================================
\documentclass[11pt,a4paper]{article}

\usepackage[utf8]{inputenc}
\usepackage[T1]{fontenc}
\usepackage[portuguese]{babel}
\usepackage[margin=2.5cm]{geometry}
\usepackage{graphicx}
\usepackage{booktabs}
\usepackage{tabularx}
\usepackage{float}
\usepackage{amsmath}
\usepackage{siunitx}
\usepackage{caption}
\usepackage[table]{xcolor}
\usepackage{listings}
\usepackage{enumitem}
\usepackage{titlesec}
\usepackage[most]{tcolorbox}
\usepackage[colorlinks=true, linkcolor=azul, citecolor=verde!85!black,
            urlcolor=azul]{hyperref}

\sisetup{output-decimal-marker={,}, per-mode=symbol}

% ----------------------------------------------------------------------------
% Paleta: verde de relvado para a estrutura, azul para o detalhe, e duas cores
% de sinal usadas só na Tarefa 3 (o que correu bem / mal).
% ----------------------------------------------------------------------------
\definecolor{verde}{HTML}{1B7F3B}
\definecolor{azul}{HTML}{1F3A5F}
\definecolor{bom}{HTML}{1B7F3B}
\definecolor{mau}{HTML}{B03A2E}
\colorlet{tabcab}{verde!16}      % linha de cabeçalho das tabelas
\colorlet{tabgrupo}{azul!8}      % linhas que separam grupos dentro das tabelas

\setlist{nosep, leftmargin=*}
\setlist[itemize]{label=\textcolor{verde}{\textbullet}}

\titleformat{\section}{\Large\bfseries\color{verde}}{\thesection}{0.6em}{}
            [{\color{verde}\titlerule[0.8pt]}]
\titleformat{\subsection}{\large\bfseries\color{azul}}{\thesubsection}{0.6em}{}
\titleformat{\paragraph}[runin]{\normalsize\bfseries\color{azul}}{}{0pt}{}

\captionsetup{labelfont={bf,color=azul}, font=small}

\newtcolorbox{caixaresumo}{colback=azul!4, colframe=azul!45, boxrule=0.5pt,
  arc=2pt, left=8pt, right=8pt, top=6pt, bottom=6pt, parbox=false}
\newtcolorbox{destaque}{colback=verde!7, colframe=verde, boxrule=0pt,
  leftrule=3pt, arc=0pt, left=8pt, right=8pt, top=5pt, bottom=5pt}

\renewcommand{\topfraction}{0.9}
\renewcommand{\textfraction}{0.07}
\renewcommand{\floatpagefraction}{0.8}

\lstset{basicstyle=\ttfamily\footnotesize, breaklines=true,
        frame=single, framesep=4pt, rulecolor=\color{azul!45},
        backgroundcolor=\color{azul!4}, keywordstyle=\color{azul}\bfseries,
        stringstyle=\color{verde!85!black},
        columns=flexible, keepspaces=true}
\renewcommand{\lstlistingname}{List}

% \fig{ficheiro}{largura}: inclui a figura se existir; senão deixa um aviso
\newcommand{\fig}[2]{\IfFileExists{#1}{\includegraphics[width=#2]{#1}}%
  {\fbox{\parbox{#2}{\centering\ttfamily\small figura em falta: \detokenize{#1}}}}}

\newcolumntype{L}{>{\raggedright\arraybackslash}X}

\title{{\color{verde}\textbf{Golf swing in MuJoCo}}\\[2mm]
       {\large\color{azul} Robótica --- Laboratório 1}}
\author{{Afonso Cruz (102815)} \and {Eduardo Batista (99931)} \and {Henrique Póvoa (102109)}}
\date{Grupo {nº\textbf{1}} \quad --- \quad 9/10/2026}

\begin{document}
\maketitle

\begin{abstract}
\begin{caixaresumo}
Simula-se em MuJoCo~\cite{todorov2012} o swing de um jogador de golfe com um
ferro~7, tratando jogador mais taco como um robot articulado. O modelo tem três
juntas atuadas (3~DoF) --- tronco, braço esquerdo e pulso --- e um braço
direito passivo preso ao punho do taco, e é controlado por binário calculado.

O swing atinge \SI{28.2}{\meter\per\second} na cabeça do taco e
\SI{33.9}{\meter\per\second} na bola, com binários dentro de limites humanos e
a sequência de um swing real. Com ruído de binário, o pulso é a junta que mais
afeta o alcance e o ombro a que mais afeta a direção. Usou-se o MuJoCo 3.13
com Python~3.11.
\end{caixaresumo}
\end{abstract}


% =============================================================================
\section{Tarefa 1 --- Modelo do jogador e do taco}
% =============================================================================

O modelo (\texttt{models/golfer\_v17.xml}) representa um adulto destro de
\SI{1.78}{\meter} e \SI{82}{\kilo\gram}, na postura de endereço com um ferro~7.
Pernas e pélvis estão fixas ao mundo; o resto move-se com os graus de liberdade
da Tabela~\ref{tab:dof}. O ficheiro XML está comentado junta a junta e não
precisa de ficheiros adicionais: todas as formas são primitivas, sem STL.

\begin{table}[H]
\centering
\caption{Graus de liberdade do modelo.}
\label{tab:dof}
\small
\begin{tabularx}{\textwidth}{@{}lLcc@{}}
\toprule
\rowcolor{tabcab}
Junta & Movimento & Amplitude & Binário máx. \\
\midrule
\rowcolor{tabgrupo}\multicolumn{4}{@{}l}{\emph{Atuadas}} \\
\texttt{torso}    & rotação do tronco em torno da coluna (inclinada \ang{35}) & \ang{-100} a \ang{95} & \SI{150}{\newton\meter} \\
\texttt{shoulder} & elevação do braço esquerdo à frente do peito, em torno da linha dos ombros & \ang{-115} a \ang{10} & \SI{100}{\newton\meter} \\
\texttt{wrist}    & taco em relação ao braço esquerdo, no plano do swing & \ang{-110} a \ang{120} & \SI{40}{\newton\meter} \\
\midrule
\rowcolor{tabgrupo}\multicolumn{4}{@{}l}{\emph{Passivas}} \\
\texttt{shoulder\_r}, \texttt{elbow\_r} & braço direito (rótula + charneira), mão presa ao punho & --- & --- \\
\texttt{neck}     & cabeça, acoplada a $-$\texttt{torso} & --- & --- \\
\texttt{ball\_free} & bola (6 dof) & --- & --- \\
\bottomrule
\end{tabularx}
\end{table}

\subsection{Juntas}
\label{sec:passivas}

\paragraph{Atuadas.}
O \texttt{torso} é a fonte de força: num corpo real ela vem das pernas e das
ancas, aqui está toda neste motor, daí os \SI{150}{\newton\meter} (um ombro
isolado produz 60--85\,\si{\newton\meter}~\cite{harbo2012}). O
\texttt{shoulder} levanta o braço esquerdo, sempre esticado, pela frente do
peito. O \texttt{wrist} roda o taco no plano do swing; os
\SI{40}{\newton\meter} são da ordem do binário que as mãos aplicam ao taco num
swing medido ($\approx$\SI{31}{\newton\meter}~\cite{nesbit2005}). Braço e taco
formam o pêndulo duplo clássico do swing~\cite{cochran1968,jorgensen1970}: o
pulso arma no topo e só abre perto do impacto~\cite{pickering1999}.

\paragraph{Passivas.}
O braço direito não tem motores. A mão está presa ao punho por uma restrição
\texttt{connect} e o braço é arrastado pelo taco: dobra no backswing e estica
no impacto, com um erro máximo na restrição de \SI{2.2}{\milli\meter}. Como o
eixo do ombro esquerdo passa pelos dois ombros, a distância que o braço direito
tem de alcançar só depende do pulso (0,54 a \SI{0.60}{\meter}, abaixo dos
\SI{0.62}{\meter} do braço esticado). A cabeça roda ao contrário do tronco
(\texttt{neck} $= -$\texttt{torso}) e fica parada a olhar para a bola. Nenhuma
junta passiva entra no modelo do controlador.

\subsection{Referenciais e convenção de sinais}

O referencial do mundo tem origem na bola, $x$ na direção do alvo e $z$ para
cima; o jogador está do lado $+y$. Na pose de endereço os eixos de todos os
corpos coincidem com os do mundo (Figura~\ref{fig:frames}); as origens estão na
Tabela~\ref{tab:frames}.

\begin{table}[H]
\centering
\caption{Origem do referencial de cada corpo.}
\label{tab:frames}
\small
\begin{tabular}{ll}
\toprule
\rowcolor{tabcab}
Corpo & Origem \\
\midrule
\texttt{torso} & anca \\
\texttt{head}  & topo do peito, sobre o eixo da coluna \\
\texttt{arm\_l} & ombro esquerdo \\
\texttt{club}  & punho (mão esquerda) \\
\texttt{upperarm\_r}, \texttt{forearm\_r} & ombro e cotovelo direitos \\
\bottomrule
\end{tabular}
\end{table}

Designa-se por $q$ o ângulo de cada junta, medido a partir da pose de endereço
($q=0$). No tronco e no pulso, $q<0$ no backswing e $q>0$ no follow-through; no
ombro, $q<0$ sempre que o braço está levantado.

\begin{destaque}
\textbf{Condição que domina todo o controlo:} a cabeça do taco só está na
posição da bola quando as três juntas valem $q=0$ ao mesmo tempo.
\end{destaque}

\begin{figure}[H]
\centering
\fig{Photos/v17_referenciais.png}{0.9\textwidth}
\caption{Referenciais dos corpos na pose de endereço ($q=0$): vermelho $x$,
verde $y$, azul $z$.}
\label{fig:frames}
\end{figure}

\subsection{Parâmetros físicos}

As massas seguem tabelas antropométricas aproximadas e as inércias são
calculadas pelo MuJoCo a partir da forma e da massa de cada \texttt{geom}
(Tabelas~\ref{tab:inercias} e~\ref{tab:params}).

\begin{table}[H]
\centering
\caption{Massa e momentos principais de inércia (no centro de massa) de cada
corpo móvel.}
\label{tab:inercias}
\small
\begin{tabular}{lcccc}
\toprule
\rowcolor{tabcab}
Corpo & Massa [\si{\kilo\gram}] & $I_1$ & $I_2$ & $I_3$ [\si{\kilo\gram\meter\squared}] \\
\midrule
\texttt{torso} (tronco + ombros)     & 28,0 & 1,020 & 0,995 & 0,217 \\
\texttt{head}                        & 6,0  & 0,032 & 0,032 & 0,021 \\
\texttt{arm\_l} (braço esq. + mão)   & 4,0  & 0,172 & 0,172 & 0,003 \\
\texttt{club}                        & 0,43 & 0,037 & 0,037 & 0,0001 \\
\texttt{upperarm\_r}                 & 2,2  & 0,028 & 0,028 & 0,002 \\
\texttt{forearm\_r} (+ mão)          & 1,8  & 0,022 & 0,022 & 0,001 \\
\texttt{ball}                        & 0,046 & \multicolumn{3}{c}{$8{,}3\times10^{-6}$} \\
\bottomrule
\end{tabular}
\end{table}

\begin{table}[H]
\centering
\caption{Restantes parâmetros físicos.}
\label{tab:params}
\small
\begin{tabularx}{\textwidth}{@{}lL@{}}
\toprule
\rowcolor{tabcab}
Parâmetro & Valor \\
\midrule
Taco & \SI{0.43}{\kilo\gram}, dos quais \SI{0.26}{\kilo\gram} na cabeça;
       \emph{loft} de \ang{32}; haste rígida~\cite{milne1992} \\
Bola & \SI{45.9}{\gram}, \SI{42.7}{\milli\meter} de diâmetro,
       \SI{6}{\milli\meter} à frente da face no endereço \\
Atrito nas juntas & amortecimento viscoso de 2, 0{,}5 e
       \SI{0.1}{\newton\meter\second} no tronco, ombro e pulso; 0{,}5 nas
       juntas do braço direito \\
Inércia de rotor & \SI{0.01}{\kilo\gram\meter\squared} em cada junta \\
Ar & densidade \SI{1.204}{\kilo\gram\per\meter\cubed}, viscosidade
       \SI{1.81e-5}{\pascal\second} \\
\bottomrule
\end{tabularx}
\end{table}

\subsection{Contacto e passo de integração}
\label{sec:contacto}

Só a cabeça do taco e a bola colidem; o corpo do jogador é apenas visual. O
impacto é definido num par de contacto dedicado (Listagem~\ref{lst:pair}).

\begin{lstlisting}[language=XML, caption={Par de contacto taco--bola em
  \texttt{golfer\_v17.xml}.}, label={lst:pair}, captionpos=b]
<contact>
  <pair geom1    = "clubhead"
        geom2    = "ball_geom"
        condim   = "3"
        solref   = "0.0004 0.15"
        solimp   = "0.95 0.99 0.001"
        friction = "0.4 0.4 0.005 0.0001 0.0001"/>
</contact>
\end{lstlisting}

\begin{itemize}
  \item \texttt{condim}: com 3, o contacto tem força normal e atrito
        tangencial.
  \item \texttt{solref}: constante de tempo de \SI{0.4}{\milli\second}, a
        duração real de um impacto de golfe~\cite{penner2003}, e razão de
        amortecimento de 0,15. Sub-amortecido, o contacto faz a bola ressaltar:
        é assim que se impõe a restituição no
        MuJoCo~\cite{todorov2012}, calibrada para \emph{smash factor}
        $\approx 1{,}2$.
  \item \texttt{solimp}: contacto quase rígido (rigidez de 0,95 a 0,99 em
        \SI{1}{\milli\meter} de penetração).
  \item \texttt{friction}: atrito tangencial de 0,4.
\end{itemize}

O passo de integração é \SI{5e-5}{\second} e desce para \SI{5e-6}{\second}
quando a cabeça do taco está a menos de \SI{15}{\centi\meter} da bola. Com o
passo normal, a velocidade da bola variava $\pm5\,\%$ conforme o instante em
que o contacto começava; com o passo fino varia $\sim$1\,\%.

\subsection{Simplificações e impacto no realismo}

\begin{table}[H]
\centering
\caption{Simplificações do modelo, consequência de cada uma e resultado que
mais afeta. A negrito, a que mais pesa nos resultados.}
\label{tab:simpl}
\small
\renewcommand{\arraystretch}{1.25}
\begin{tabularx}{\textwidth}{@{}>{\raggedright\arraybackslash}p{3.5cm}L>{\raggedright\arraybackslash}p{2.7cm}@{}}
\toprule
\rowcolor{tabcab}
Simplificação & Consequência & Afeta \\
\midrule
\rowcolor{tabgrupo}\multicolumn{3}{@{}l}{\textbf{Corpo}} \\
\textbf{Pernas e pélvis rígidas} & Sem rotação das ancas nem transferência de peso~\cite{hume2005} & \textbf{velocidade} \\
Cotovelo esq.\ rígido  & O braço da frente não dobra depois do impacto        & follow-through \\
Ombro esq.\ com 1 dof  & Um ombro real tem 3 e a omoplata desliza             & plano do swing \\
Pulso com 1 dof        & A face não abre nem fecha: sem \emph{slice} nem \emph{hook} & direção \\
Braço direito passivo  & É só carga, não empurra o taco                       & velocidade \\
Corpo sem contactos    & Os braços podem atravessar o corpo (medido: \SI{4.9}{\milli\meter}) & aspeto visual \\
\addlinespace
\rowcolor{tabgrupo}\multicolumn{3}{@{}l}{\textbf{Taco, impacto e voo}} \\
Taco rígido            & Sem flexão da haste~\cite{milne1992} nem geometria real da face & ponto de impacto \\
Contacto mole          & Restituição calibrada; a bola sai pela normal da face & lançamento e \emph{spin} \\
$C_d$ e $C_L$ constantes & Magnus linear na rotação, sem saturação             & voo \\
Relva rígida           & Sem travagem pelo backspin; usa-se o \emph{carry}~\cite{penner2002} & rolamento \\
\addlinespace
\rowcolor{tabgrupo}\multicolumn{3}{@{}l}{\textbf{Atuação e controlo}} \\
Motores ideais         & Sem relação força--velocidade nem atraso de ativação & binário disponível \\
Controlador exato      & Um humano tem atrasos de $\sim$\SI{100}{\milli\second} e um modelo aproximado & rejeição do ruído \\
\bottomrule
\end{tabularx}
\end{table}
% =============================================================================
\section{Tarefa 2 --- Controlo e ensaios}
% =============================================================================

\subsection{Trajetória e controlador}
\label{sec:controlo}

Cada junta segue uma referência de polinómios de 5.º grau, contínua em posição,
velocidade e aceleração~\cite{sequeira2024}: backswing de \SI{0.80}{\second}
até ao topo $(\ang{-90}, \ang{-85}, \ang{-95})$, pausa de \SI{0.05}{\second},
downswing de \SI{0.30}{\second} e follow-through até
$(\ang{85}, \ang{-70}, \ang{105})$. Como num swing humano, no downswing o
tronco arranca primeiro, o braço a \SI{10}{\percent} e o pulso fica armado até
\SI{55}{\percent}. As três juntas chegam a $q=0$ no mesmo instante
($t=\SI{1.15}{\second}$): o braço parado, o tronco ainda a rodar e o pulso no
pico de velocidade.

A referência é seguida por binário calculado~\cite{siciliano2009,sequeira2024},
\begin{equation}
  \boldsymbol\tau = M(q)\left[\ddot q_{\text{ref}}
      + K_d\left(\dot q_{\text{ref}} - \dot q\right)
      + K_p\left(q_{\text{ref}} - q\right)\right] + h(q,\dot q),
  \label{eq:ct}
\end{equation}
com $M$ a matriz de massa das três juntas atuadas, $h$ os termos de Coriolis,
centrífugos e de gravidade, $K_p = 900$ e $K_d = 60$. Os binários são saturados
nos limites do XML.

\subsection{Bateria de testes}

\begin{table}[H]
\centering
\caption{Bateria de testes, pela ordem em que foram corridos.}
\label{tab:testes}
\small
\begin{tabular}{@{}rlll@{}}
\toprule
\rowcolor{tabcab}
\# & Teste & O que verifica & Resultado \\
\midrule
1 & Pêndulo simples e pêndulo com bola  & física de base                  & erro do período $<0{,}2\,\%$ \\
2 & Convergência do passo de integração & repetibilidade do impacto       & Secção~\ref{sec:contacto} \\
3 & Bola em voo livre                   & $C_d$ e $C_L$ efetivos          & Tabela~\ref{tab:aero} \\
4 & Swing nominal                       & comparação com um ferro~7 real  & Tabela~\ref{tab:nominal} \\
5 & Sequência cinemática                & ordem dos picos de velocidade   & Figura~\ref{fig:seq} \\
6 & Ruído de binário em cada junta      & sensibilidade a perturbações    & Secção~\ref{sec:pert} \\
7 & Três modelos de ar no corpo         & arrasto no braço e no taco      & Secção~\ref{sec:aero} \\
8 & ``Tacada perfeita''                 & voo com um lançamento real      & Secção~\ref{sec:aero} \\
\bottomrule
\end{tabular}
\end{table}

Os testes correm em Python (\texttt{src/}, instruções no \texttt{README.md}).
O XML abre-se também na \emph{Simulate App} do MuJoCo, onde os painéis
\emph{Joint} e \emph{Sensor} permitem mover cada junta à mão e ver ângulos,
velocidades e binários ao vivo.

\subsection{Ensaio nominal}
\label{sec:nominal}

\begin{figure}[H]
\centering
\fig{Photos/v17_sequencia_face_on.png}{\textwidth}\\[1mm]
\fig{Photos/v17_sequencia_down_the_line.png}{\textwidth}
\caption{Fases do swing simulado, câmaras \emph{face-on} (cima) e
\emph{down-the-line} (baixo). O braço direito dobra no topo e estica no impacto
sem ser comandado, e a cabeça fica parada.}
\label{fig:fases}
\end{figure}

\begin{figure}[H]
\centering
\fig{Photos/v17_sequencia.png}{0.7\textwidth}
\caption{Ângulos, velocidades e binários das três juntas atuadas. O pulso
atinge a velocidade máxima no impacto, depois de tronco e ombro travarem.}
\label{fig:seq}
\end{figure}

A Figura~\ref{fig:seq} confirma três coisas. A referência é seguida quase sem
erro: as juntas cruzam zero no impacto com desvios de $-1{,}0$, $+1{,}1$ e
$+2{,}7$\,\si{\degree}. A velocidade passa dos segmentos grandes para o taco,
como num swing real: ombro e tronco atingem o pico (598 e
\SI{453}{\degree\per\second}) antes do impacto e travam, e o pulso chega a
\SI{1392}{\degree\per\second} exatamente no impacto. E os binários ficam dentro
de limites humanos: tronco e ombro nunca saturam e o pulso satura em
\SI{13}{\percent} do downswing.

Face a um ferro~7 real (Tabela~\ref{tab:nominal}), o tempo e a sincronização
estão certos; a velocidade, a rotação e o alcance ficam abaixo e o lançamento
sai alto.

\begin{table}[H]
\centering
\caption{Swing nominal e médias publicadas de um ferro~7 (TrackMan, LPGA a PGA
Tour~\cite{trackman2024}; amador: \cite{golfmonthly_7iron}).}
\label{tab:nominal}
\small
\begin{tabular}{lcc}
\toprule
\rowcolor{tabcab}
Grandeza & Modelo & Real \\
\midrule
Duração do downswing          & \SI{0.30}{\second} & 0,26--0,39~\cite{grober2010} \\
Velocidade da cabeça do taco  & \SI{28.2}{\meter\per\second} & 34,0--40,2 (amador $\sim$36) \\
Velocidade da bola            & \SI{33.9}{\meter\per\second} & 44,7--53,6 \\
\emph{Smash factor}           & 1,20 & 1,33--1,38 \\
Ângulo de ataque              & \ang{-1.6} & $\sim\ang{-4}$ \\
Ângulo de lançamento          & \ang{33.6} & 16--19 \\
\emph{Spin}                   & \SI{2510}{rpm} & 6700--7100 \\
\emph{Carry}                  & \SI{94.6}{\meter} & 129--157 \\
Altura máxima                 & \SI{21.1}{\meter} & 24--29 \\
Ângulo de descida             & \ang{44.4} & 47--50 \\
\bottomrule
\end{tabular}
\end{table}

\subsection{Aerodinâmica}
\label{sec:aero}

O ar está ligado para todo o modelo. A bola tem arrasto e sustentação de
Magnus~\cite{mujoco_fluid}, e os coeficientes efetivos, medidos em voo livre,
ficam dentro das gamas publicadas (Tabela~\ref{tab:aero}).

\begin{table}[H]
\centering
\caption{Coeficientes aerodinâmicos efetivos da bola, medidos em voo livre nas
condições de lançamento do swing.}
\label{tab:aero}
\small
\begin{tabular}{lcc}
\toprule
\rowcolor{tabcab}
 & Modelo & Literatura \\
\midrule
$C_d$ & 0,28 & 0,27--0,32~\cite{bearman1976,smits1994} \\
$C_L$ & 0,22 & 0,15--0,25~\cite{bearman1976,smits1994} \\
Decaimento da rotação & \SI{2.8}{\percent\per\second} & 4--5\,\si{\percent\per\second}~\cite{smits1994,tavares1999} \\
\bottomrule
\end{tabular}
\end{table}

\begin{figure}[H]
\centering
\fig{Photos/v17_aerodinamica.png}{\textwidth}
\caption{Esquerda: voo da bola no vácuo, só com arrasto e com
arrasto\,+\,Magnus. Direita: velocidade da cabeça do taco com três modelos de
ar no braço e no taco.}
\label{fig:aero}
\end{figure}

\paragraph{Voo.}
A bola lançada sozinha com as condições de saída do swing aterra a
\SI{94.6}{\meter}, o mesmo valor do swing completo; sem Magnus daria
\SI{78}{\meter} e no vácuo \SI{108}{\meter}.

\paragraph{Ar no corpo.}
O modelo de ar por omissão do MuJoCo dá \SI{11.8}{\newton\meter} de arrasto no
pulso, seis vezes mais do que o modelo por elipsoides
(\SI{1.8}{\newton\meter}) e do que uma estimativa à
mão~\cite{hoerner1965} (\SI{1.5}{\newton\meter}); usam-se por isso elipsoides.
Com eles, o ar não altera a velocidade do taco.

\paragraph{Limitação.}
Lançada com as condições médias do LPGA Tour (\SI{46.5}{\meter\per\second},
\ang{19}, \SI{6700}{rpm}), a bola voa \SI{162}{\meter} com \SI{44}{\meter} de
altura, contra 129 e \SI{24}{\meter} publicados. O Magnus do MuJoCo é linear na
rotação, enquanto numa bola real satura~\cite{smits1994}.

\subsection{Ensaios com perturbações}
\label{sec:pert}

Somou-se ruído gaussiano de desvio padrão $\sigma = \SI{10}{\newton\meter}$ ao
binário de uma junta atuada de cada vez e repetiu-se o swing $N=20$ vezes por
junta. As juntas passivas não têm motor onde somar ruído; são perturbadas
indiretamente, através do taco e do tronco. O ruído é mantido constante
durante \SI{0.1}{\second} (\SI{10}{\hertz}, a banda do tremor
fisiológico~\cite{elble1990,deuschl1998}) e aplicado ao binário porque o ruído
motor humano cresce com o comando~\cite{harris1998}.

Na Tabela~\ref{tab:pert}, $\sigma/\tau_{\max}$ é a fração do motor que o ruído
representa, ``Válidos'' são os swings em que o taco acertou na bola e
``Lateral'' é o desvio no ponto de aterragem.

\begin{table}[H]
\centering
\caption{Efeito do ruído de binário, média $\pm$ desvio padrão dos swings
válidos. Lateral: $+$ = esquerda do alvo, $-$ = direita.}
\label{tab:pert}
\small
\begin{tabular}{lcccccc}
\toprule
\rowcolor{tabcab}
Junta & $\sigma/\tau_{\max}$ & Válidos & $v_{\text{cab}}$ [\si{\meter\per\second}] & Lanç. [\si{\degree}] & \emph{Carry} [\si{\meter}] & Lateral [\si{\meter}] \\
\midrule
\texttt{torso}    & 7\,\%  & 20/20 & $28{,}2 \pm 0{,}0$ & $33{,}8 \pm 0{,}2$ & $95{,}3 \pm 0{,}4$ & $-0{,}8 \pm 0{,}2$ \\
\texttt{shoulder} & 10\,\% & 19/20 & $28{,}2 \pm 0{,}1$ & $33{,}7 \pm 0{,}1$ & $95{,}4 \pm 0{,}6$ & $-1{,}4 \pm 2{,}6$ \\
\texttt{wrist}    & 25\,\% & 20/20 & $28{,}1 \pm 0{,}4$ & $33{,}6 \pm 0{,}8$ & $94{,}7 \pm 2{,}3$ & $-0{,}5 \pm 0{,}8$ \\
\bottomrule
\end{tabular}
\end{table}

\begin{figure}[H]
\centering
\fig{Photos/v17_dispersao.png}{\textwidth}
\caption{Pontos de aterragem dos swings com ruído, um painel por junta
(estrela: nominal; barras: média $\pm$ 1 desvio padrão).}
\label{fig:dispersao}
\end{figure}

\begin{itemize}
  \item \textbf{Tronco:} quase insensível; o \emph{carry} varia apenas
        $\pm$\SI{0.4}{\meter}.
  \item \textbf{Pulso:} o mais sensível no alcance. É a única junta em que o
        ruído se vê na velocidade ($\pm$\SI{0.4}{\meter\per\second}) e no
        \emph{carry} ($\pm$\SI{2.3}{\meter}), porque \SI{10}{\newton\meter}
        são \SI{25}{\percent} do seu motor.
  \item \textbf{Ombro:} o mais sensível na direção. Dois swings aterraram 5 e
        \SI{11}{\meter} à direita e um falhou a bola: um erro de \ang{0.7}
        desloca a cabeça do taco quase \SI{1}{\centi\meter} na vertical e a
        bola é apanhada fora do centro da face.
\end{itemize}

\newpage
% =============================================================================
\section{Tarefa 3 --- Comentário aos resultados}
\label{sec:comentario}
% =============================================================================

Com apenas três graus de liberdade atuados, o modelo executa um swing completo
e anatomicamente plausível. O movimento e o controlo saíram bem; os números do
impacto e do voo ficaram aquém de um ferro~7 real, por razões que se conseguem
identificar e corrigir.

\subsection*{\textcolor{bom}{Pontos bons}}

\begin{itemize}[itemsep=3pt]
  \item \textbf{Três juntas chegam para um swing real.} O downswing dura
        \SI{0.30}{\second} (real: 0,26--0,39). Tronco e ombro atingem o pico
        de velocidade antes do impacto e travam, e o pulso chega ao máximo
        (\SI{1392}{\degree\per\second}) exatamente no impacto, como prevê o
        pêndulo duplo~\cite{jorgensen1970,pickering1999}.
  \item \textbf{Impacto preciso, com binários humanos.} As três juntas cruzam
        a pose de endereço com desvios de $-1{,}0$, $+1{,}1$ e \ang{+2.7}.
        Tronco e ombro nunca saturam; o pulso satura em \SI{13}{\percent} do
        downswing.
  \item \textbf{Movimento plausível.} O braço direito, sem motores, acompanha
        o taco com um erro máximo de \SI{2.2}{\milli\meter}, a cabeça fica
        parada e nenhuma parte atravessa o corpo em mais de
        \SI{4.9}{\milli\meter}.
  \item \textbf{Robusto e fácil de interpretar.} Com ruído de
        \SI{10}{\newton\meter} em cada junta, 59 de 60 swings acertam na bola
        e o \emph{carry} varia no máximo $\pm$\SI{2.3}{\meter}. Com poucas
        juntas, o papel de cada uma fica claro: o pulso afeta o alcance, o
        ombro a direção e o tronco quase nada.
  \item \textbf{Voo validado.} $C_d$ e $C_L$ ficam dentro das gamas
        publicadas.
\end{itemize}

\subsection*{\textcolor{mau}{O que não correu tão bem}}

\begin{itemize}[itemsep=3pt]
  \item \textbf{Velocidade baixa.} A cabeça do taco chega a
        \SI{28.2}{\meter\per\second}, abaixo de um amador (34--40). Com as
        ancas fixas, só o tronco e o pulso contribuem.
  \item \textbf{Lançamento alto e pouca rotação.} A bola sai a \ang{33.6}
        (real: 16--\ang{19}) com \SI{2510}{rpm} (real: 6700--7100), porque
        sai pela normal da face. O \emph{smash factor} de 1,20 (real: 1,33)
        também é um parâmetro calibrado e não uma propriedade do material.
  \item \textbf{Alcance curto.} O \emph{carry} é de \SI{94.6}{\meter} (real:
        129--157), consequência dos dois pontos anteriores. A altura do voo só
        se aproxima da real por compensação entre lançamento alto e rotação
        baixa.
  \item \textbf{Sem erros de face e dispersão otimista.} A face não abre nem
        fecha, logo o modelo não reproduz \emph{slice} nem \emph{hook}. E o
        controlador conhece a dinâmica exata, por isso rejeita o ruído melhor
        do que um humano.
\end{itemize}

\subsection*{\textcolor{azul}{O que podíamos melhorar}}

\begin{itemize}[itemsep=3pt]
  \item \textbf{Libertar as ancas}, para ganhar a rotação e a transferência de
        peso que faltam à velocidade.
  \item \textbf{Inclinar a haste no impacto}, com mais um grau de liberdade,
        para baixar o lançamento e aumentar a rotação.
  \item \textbf{Acrescentar a rotação dos antebraços}, para a face abrir e
        fechar e o ensaio de ruído mostrar erros de direção realistas.
  \item \textbf{Dar atrasos e um modelo aproximado ao controlador}, para a
        dispersão com ruído se aproximar da de um jogador.
\end{itemize}
\newpage
% =============================================================================
\section*{Contribuição dos membros}
Todos os membros contribuíram de forma igual para a realização do projeto.
Cada um ficou responsável por uma área, sem deixar de participar nas restantes:
\begin{itemize}
  \item \textbf{Afonso Cruz:} código Python de simulação do modelo do golfista;
        redação do relatório.
  \item \textbf{Eduardo Batista:} modelo aerodinâmico (arrasto e efeito de
        Magnus) e trajetória da bola.
  \item \textbf{Henrique Póvoa:} modelo XML do golfista e do taco e respetiva
        simulação; redação do relatório.
\end{itemize}

% =============================================================================
\bibliographystyle{IEEEtran}
\bibliography{refs}

\end{document}